\documentclass[11pt,a4paper]{article}

\def\nyear {2026}
\def\nterm {Winter}
\def\nlecturer {Liran Rotem}
\def\ncourse {Functional Analysis}

\makeatletter

% packages
\usepackage{amssymb,amsfonts,amsmath,calc,tikz,pgfplots,geometry,mathtools}
\usepackage{color}   % May be necessary if you want to color links
\usepackage[svgnames]{xcolor}
\usepackage[hidelinks]{hyperref}
\usepackage{forest}
\usepackage{commath} % ew
\usepackage{esvect} % \vv makes vectors
\usepackage{centernot} % for the comparison
\usepackage{esdiff}
\usepackage{esint}
\usepackage{amsthm}
\usepackage{fancyhdr}
\usepackage{bm}
\usepackage{witharrows}
\usepackage{bookmark}
\usepackage{tikz-cd}
\usepackage{quiver}
\usepackage{bbm}
\usepackage{textcomp}
\usepackage{float}
\usepackage{gensymb}
\usepackage{cleveref}
\usepackage{environ}
\usepackage{comment}
\usepackage{cancel}
\usepackage{xparse}

% Inevitable cs
\usepackage{listings}
\usepackage[]{algorithm2e}

% tikz libraries
\usetikzlibrary{positioning}
\usetikzlibrary{matrix}
\usetikzlibrary{arrows}
\usetikzlibrary{bending}
\usetikzlibrary{arrows.meta}
\usetikzlibrary{decorations.markings}
\usetikzlibrary{decorations.pathreplacing}
\usetikzlibrary{decorations.markings}
\usetikzlibrary{patterns}
\usetikzlibrary{backgrounds}

% Page style setup
\pagestyle{fancy}
\fancyhfoffset{0pt}
\renewcommand{\headwidth}{\textwidth}
\geometry{margin=1in}
\pgfplotsset{compat=1.18}
\setlength{\headheight}{14.6pt}
\addtolength{\topmargin}{-1.6pt}
\hypersetup{
    colorlinks=false,
    linktoc=section,
    linkcolor=black,
}

%% maketitle setup
\ifx \nauthor\undefined
  \def\nauthor{Yeheli Fomberg}
\else
\fi

\ifx \nemail\undefined
  \def\nemail{\href{mailto:yp@campus.technion.ac.il}
  {\texttt{<yp@campus.technion.ac.il>}}}
\else
\fi

\ifx \ncoursehead \undefined
  \def\ncoursehead{\ncourse}
\fi

\lhead{\emph{\nouppercase{\leftmark}}}
\ifx \nextra \undefined
  \rhead{
    \ifnum\thepage=1
    \else
      \ncoursehead
    \fi}
\else
  \rhead{
    \ifnum\thepage=1
    \else
      \ncoursehead \ (\nextra)
    \fi}
\fi

\def\notes{notes}
\def\problems{problems}

\ifx \ntype \undefined
  \def\ntype{notes}
\else
\fi

\ifx \ntype \notes
  \let\@real@maketitle\maketitle
  \renewcommand{\maketitle}{\@real@maketitle\begin{center}
  \begin{minipage}[c]{0.9\textwidth}\centering\footnotesize
  These notes are not endorsed by the lecturers.
  I have revised them outside lectures to incorporate supplementary
  explanations, clarifications, and material for fun.
  While I have strived for accuracy, any errors or misinterpretations 
  are most likely mine.
  \end{minipage}\end{center}}
\else
  \ifx \ntype \problems
    \let\@real@maketitle\maketitle
    \renewcommand{\maketitle}{\@real@maketitle\begin{center}
    \begin{minipage}[c]{0.9\textwidth}\centering\footnotesize
    These problems are mostly based on the problem sets I was given
    when taking the course.
    
    While I have strived for accuracy and completeness,
    some of the problems would not have been here without the help of
    friends.
    \end{minipage}\end{center}}
  \else
  \fi
\fi


% theorem environments
\theoremstyle{definition}
\newtheorem{definition}{Definition}[section]
\newtheorem{remark}{Remark}[section]
\newtheorem{example}{Example}[section]
\newtheorem{exercise}{Exercise}[section]
\newtheorem{paradox}{Paradox}[section]
\newtheorem{entry}{Entry}[section]
\newtheorem*{solution}{Solution}
\newtheorem*{question}{Question}
\newtheorem*{answer}{Answer}
\newtheorem*{problem}{Problem}
\theoremstyle{plain}
\newtheorem{theorem}{Theorem}[section]
\newtheorem{proposition}[theorem]{Proposition}
\newtheorem{lemma}[theorem]{Lemma}
\newtheorem{corollary}[theorem]{Corollary}
\newenvironment{proofof}[1]{%
    \renewcommand{\proofname}{Proof of #1}%
    \begin{proof}
}{%
    \end{proof}
}
% tikz customization
\pgfarrowsdeclarecombine{twolatex'}{twolatex'}{latex'}{latex'}{latex'}{latex'}
\tikzset{->/.style = {decoration={markings,
                                  mark=at position 0.999
                                  with {\arrow[scale=2]{latex'}}},
                      postaction={decorate}}}
\tikzset{<-/.style = {decoration={markings,
                                  mark=at position 0.001 with {\arrowreversed[scale=2]{latex'}}},
                      postaction={decorate}}}
\tikzset{-->/.style = {decoration={markings,
                                  mark=at position 0.999
                                  with {\arrow[scale=2]{latex'[sep=0pt -2.5]}}},
                      postaction={decorate}}}
\tikzset{<--/.style = {decoration={markings,
                                  mark=at position 0.001
                                  with {\arrowreversed[scale=2]{latex'[sep=0pt -2.5]}}},
                      postaction={decorate}}}
\tikzset{<->/.style = {decoration={markings,
                                  mark=at position 0.001 with {\arrowreversed[scale=2]{latex'}},
                                  mark=at position 0.999 with {\arrow[scale=2]{latex'}},},
                      postaction={decorate}}}
\tikzset{->-/.style = {decoration={markings,
    mark=at position 0.50+0.225/#1 with {\arrow[scale=2]{latex'}}},
                      postaction={decorate}}}
\tikzset{-<-/.style = {decoration={markings,
                                   mark=at position 0.50-0.225/#1 with {\arrowreversed[scale=2]{latex'}}},
                       postaction={decorate}}}
\tikzset{->>/.style = {decoration={markings,
                                  mark=at position 1 with {\arrow[scale=2]{latex'}}},
                      postaction={decorate}}}
\tikzset{<<-/.style = {decoration={markings,
                                  mark=at position 0 with {\arrowreversed[scale=2]{twolatex'}}},
                      postaction={decorate}}}
\tikzset{<<->>/.style = {decoration={markings,
                                   mark=at position 0 with {\arrowreversed[scale=2]{twolatex'}},
                                   mark=at position 1 with {\arrow[scale=2]{twolatex'}}},
                       postaction={decorate}}}
\tikzset{->>-/.style = {decoration={markings,
                                   mark=at position 0.50+0.450/#1 with {\arrow[scale=2]{twolatex'}}},
                       postaction={decorate}}}
\tikzset{-<<-/.style = {decoration={markings,
                                   mark=at position 0.50-0.450/#1 with {\arrowreversed[scale=2]{twolatex'}}},
                       postaction={decorate}}}

\pgfarrowsdeclare{biggertip}{biggertip}{%
  \setlength{\arrowsize}{1pt}
  \addtolength{\arrowsize}{.1\pgflinewidth}
  \pgfarrowsrightextend{0}
  \pgfarrowsleftextend{-5\arrowsize}
}{%
  \setlength{\arrowsize}{1pt}
  \addtolength{\arrowsize}{.1\pgflinewidth}
  \pgfpathmoveto{\pgfpoint{-5\arrowsize}{4\arrowsize}}
  \pgfpathlineto{\pgfpointorigin}
  \pgfpathlineto{\pgfpoint{-5\arrowsize}{-4\arrowsize}}
  \pgfusepathqstroke
}
\tikzset{
	EdgeStyle/.style = {>=biggertip}
}

\tikzset{circ/.style = {fill, circle, inner sep = 0, minimum size = 3}}
\tikzset{scirc/.style = {fill, circle, inner sep = 0, minimum size = 1.5}}
\tikzset{mstate/.style={circle, draw, black, text=black, minimum width=0.7cm}}

\tikzset{eqpic/.style={baseline={([yshift=-.5ex]current bounding box.center)}}}

\definecolor{mblue}{rgb}{0.2, 0.3, 0.8}
\definecolor{morange}{rgb}{1, 0.5, 0}
\definecolor{mgreen}{rgb}{0, 0.4, 0.2}
\definecolor{mred}{rgb}{0.5, 0, 0}

% topology
\DeclareMathOperator{\dfr}{dfr} % deformation retract
\newcommand{\Cells}{\text{Cells}}
\newcommand{\RP}{\mathbb{RP}}
\newcommand{\CP}{\mathbb{CP}}

% order theory
\newcommand{\join}{\vee}
\newcommand{\meet}{\wedge}

% algebraic geometry
\DeclareMathOperator{\Rad}{Rad} % Radical of modules
\DeclareMathOperator{\Spec}{Spec}
% \renewcommand{\div}{\mathrm{div}} % divisors
\DeclareMathOperator{\Mer}{Mer} % Meromorphic functions
\DeclareMathOperator{\Exc}{Exc} % Exceptional divisor
\DeclareMathOperator{\Td}{Td} % Todd classes

% algebra

\DeclareMathOperator{\trdeg}{tr.deg.}
\DeclareMathOperator{\odd}{odd}
\DeclareMathOperator{\even}{even}
\DeclareMathOperator{\lcm}{lcm}
\DeclareMathOperator{\ord}{ord}
\DeclareMathOperator{\codim}{codim}
\DeclareMathOperator{\Ann}{Ann}
\DeclareMathOperator{\Ext}{Ext}
\DeclareMathOperator{\Sym}{Sym}
\DeclareMathOperator{\Out}{Out}
\DeclareMathOperator{\Aut}{Aut}
\DeclareMathOperator{\End}{End}
\DeclareMathOperator{\Inn}{Inn}
\DeclareMathOperator{\Mat}{Mat}
\DeclareMathOperator{\Fix}{Fix}
\DeclareMathOperator{\std}{std}
\DeclareMathOperator{\sgn}{sgn}
\DeclareMathOperator{\adj}{adj}
\DeclareMathOperator{\id}{id}
\DeclareMathOperator{\op}{op}
\DeclareMathOperator{\tr}{tr}
\DeclareMathOperator{\diag}{diag}
\DeclareMathOperator{\rank}{rank}
\DeclareMathOperator{\rk}{rk}
\DeclareMathOperator{\coker}{coker}
\DeclareMathOperator{\Mult}{Mult} % multiplier algebra
\DeclareMathOperator{\Frac}{Frac} % fraction field
\DeclareMathOperator{\Rat}{Rat}
\DeclareMathOperator{\Gal}{Gal}
\newcommand{\ioplus}{\overset{\text{int}}{\oplus}} % interior direct sum
\newcommand{\GG}{\mathcal G} % Galois correspondence
\newcommand{\FF}{\mathcal F} % Galois correspondence
\newcommand{\cupp}{\smile} % cup product
\newcommand{\capp}{\frown} % cap product
\newcommand{\actson}{\curvearrowright}
\newcommand{\bra}{\left\langle}
\newcommand{\ket}{\right\rangle}
\newcommand{\idealin}{\triangleleft\,}
\newcommand{\normalin}{\triangleleft\,}
\newcommand{\ip}[2]{\left\langle #1, #2 \right\rangle}
\newcommand{\bigslant}[2]
{{\raisebox{.2em}{$#1$}\left/\raisebox{-.2em}{$#2$}\right.}}
\newcommand{\properideal}{%
  \mathrel{\ooalign{$\lneq$\cr\raise.22ex\hbox{$\lhd$}\cr}}}

% categories
\newcommand{\cat}[1]{\mathbf{#1}}
\newcommand{\Ch}{\cat{Ch}} % Chain complexes   
\newcommand{\Set}{\cat{Set}}
\newcommand{\Grp}{\cat{Grp}}
\newcommand{\Ab}{\cat{Ab}}
\newcommand{\Ring}{\cat{Ring}}
\newcommand{\CRing}{\cat{CRing}}
\newcommand{\Top}{\cat{Top}}
\newcommand{\Vect}{\cat{Vect}}
\newcommand{\cmod}{\cat{mod}}

%% matrix groups
\newcommand{\GL}{\mathrm{GL}}
\newcommand{\Or}{\mathrm{O}}
\newcommand{\PGL}{\mathrm{PGL}}
\newcommand{\PSL}{\mathrm{PSL}}
\newcommand{\PSO}{\mathrm{PSO}}
\newcommand{\PSU}{\mathrm{PSU}}
\newcommand{\SL}{\mathrm{SL}}
\newcommand{\msl}{\mathrm{sl}}
\newcommand{\SO}{\mathrm{SO}}
\newcommand{\Spin}{\mathrm{Spin}}
\newcommand{\Sp}{\mathrm{Sp}}
\newcommand{\SU}{\mathrm{SU}}
\newcommand{\U}{\mathrm{U}}
\let\char\relax
\newcommand{\char}{\text{char}}

% analysis
\renewcommand{\Re}{\mathrm{Re}}
\renewcommand{\Im}{\mathrm{Im}}
\NewDocumentCommand{\dx}{o}{
  \IfNoValueTF{#1}
    {\dif x}                  % if no optional argument
    {\dif^{#1} \!x}         % if optional argument
}
\NewDocumentCommand{\dy}{o}{
  \IfNoValueTF{#1}
    {\dif y}                  % if no optional argument
    {\dif^{#1} \!y}         % if optional argument
}
\newcommand{\dt}{\dif t}
\newcommand{\du}{\dif u}
\newcommand{\dv}{\dif v}
\newcommand{\dz}{\dif z}
\newcommand{\ds}{\dif s}
\newcommand{\dw}{\dif w}
\newcommand{\dr}{\dif r}
\newcommand{\dm}{\dif m}
\newcommand{\df}{\dif f}
\newcommand{\dV}{\dif V}
\newcommand{\dS}{\dif S}
\newcommand{\dA}{\dif \mathrm{A}}
\newcommand{\dmu}{\dif \mu}
\newcommand{\dnu}{\dif \nu}
\newcommand{\dtheta}{\dif \theta}
\newcommand{\domega}{\dif \omega}
\newcommand{\dsigma}{\dif \sigma}
\newcommand{\dtau}{\dif \tau}
\newcommand{\dvarphi}{\dif \varphi}
\renewcommand{\dif}{\mathrm{d}}
\renewcommand{\Dif}{\mathrm{D}}
\renewcommand{\triangle}{\Delta}
\newcommand{\ptriangle}{\partial \triangle}
\DeclareMathOperator{\im}{im}
\DeclareMathOperator{\cis}{cis}
\DeclareMathOperator{\rot}{rot}
\DeclareMathOperator{\vspan}{span}
\DeclareMathOperator{\grad}{grad}
\DeclareMathOperator{\curl}{curl}
\DeclareMathOperator{\esssup}{esssup}
\newcommand{\hodge}{{\mathnormal{\star}}}
\DeclareMathOperator{\range}{range}
\DeclareMathOperator{\Hom}{Hom}
\DeclareMathOperator{\Int}{Int}
\DeclareMathOperator{\Conv}{Conv} % convex hull
\newcommand{\ind}{\mathbbm{1}}
\DeclareMathOperator{\Res}{Res}
\DeclareMathOperator{\graph}{graph}
\DeclareMathOperator{\diam}{diam}
\DeclareMathOperator{\supp}{supp}
\DeclareMathOperator{\Vol}{Vol} % Volume
\DeclareMathOperator{\Area}{Area} % Area
\DeclareMathOperator{\area}{area}
\DeclareMathOperator{\Ind}{Ind} % Index
\newcommand{\length}{\mathrm{length}}

% logic
\DeclareMathOperator{\MOD}{MOD}
\DeclareMathOperator{\Theory}{Theory}
\newcommand{\notimplies}{%
  \mathrel{{\ooalign{\hidewidth$\not\phantom{=}$\hidewidth\cr$\implies$}}}}

% nice
\newcommand{\half}{\frac{1}{2}}
\newcommand{\pair}{\del}
\newcommand{\taking}[1]{\xrightarrow{#1}}
\newcommand{\otaking}[1]{\xleftarrow{#1}}
\newcommand{\inv}{^{-1}}
\newcommand{\ot}{\leftarrow}
\newcommand{\ninfty}{-\infty}
\newcommand{\pinfty}{+\infty}
\newcommand{\floor}[1]{\left\lfloor #1 \right\rfloor}
\newcommand{\ceil}[1]{\left\lceil #1 \right\rceil}
\newcommand{\viff}{\Big\Updownarrow} % vertical iff
\newcommand{\bmid}{\,\middle\vert\,} % big mid

% probability
\newcommand{\Prob}{\mathbf{P}}
\renewcommand{\vec}[1]{\boldsymbol{\mathbf{#1}}}
\DeclareMathOperator{\Bin}{Bin}
\DeclareMathOperator{\Geo}{Geo}
\DeclareMathOperator{\Poi}{Poi}
\DeclareMathOperator{\Exp}{Exp}
\DeclareMathOperator{\Var}{Var} % Variance
\DeclareMathOperator{\Cov}{Cov}

% special letters
\newcommand{\N}{\mathbb{N}}
\newcommand{\Z}{\mathbb{Z}}
\newcommand{\Q}{\mathbb{Q}}
\newcommand{\R}{\mathbb{R}}
\newcommand{\C}{\mathbb{C}}
\newcommand{\F}{\mathbb{F}}
\newcommand{\E}{\mathbf{E}}
\newcommand{\D}{\mathbb{D}}
\renewcommand{\H}{\mathbb{H}}
\newcommand{\ps}{\mathcal{P}}
\newcommand{\M}{\mathcal{M}}
\renewcommand{\L}{\mathcal{L}}
\renewcommand{\O}{\mathcal{O}}
\renewcommand{\U}{\mathcal{U}}
\newcommand{\Omicron}{O}
\newcommand{\powerset}{\mathcal{P}}

% text
\newcommand{\st}{\text{ s.t. }}
\newcommand{\tand}{\quad \text{and} \quad}
\newcommand{\tor}{\quad \text{or} \quad}
\newcommand{\stand}{\text{ and }}
\newcommand{\stor}{\text{ or }}
\renewcommand{\tt}[1]{\textnormal{\textbf{(#1).}}} %tt=theorem title GET RID OF

% title format
\title{\textbf{\ncourse}}

\ifx \ntype \notes
  \author{Notes taken by \nauthor\, \nemail \\ Based on lectures by \nlecturer}
  \date{\nterm\ \nyear}
\else
  \ifx \ntype \problems
    \author{These problems were solved by \nauthor \\ \nemail}
  \else
  \fi
\fi

\makeatother


\begin{document}
\maketitle

% Insert cool image here

\newpage
\tableofcontents
\newpage

\section{Introduction}

\begin{definition}[Normed space]
  Let $V$ be a vector space over a field $\mathbb F$ (usually $\R$ or $\C$).
  A norm $\norm{\cdot} \colon V \to \R$ is a function that satisfies
  the following properties:
  \begin{enumerate}
    \item[(1)] $\norm{x} \geq 0$ and $\norm{x} = 0$ if and only if $x = 0$.
    \item[(2)] $\norm{\lambda x} = \abs{\lambda} \cdot \norm{x}$
      for every $\lambda \in \mathbb F$ and $x \in V$.
    \item[(3)] $\norm{x + y} \le \norm{x} + \norm{y}$ for all $x,y \in V$.
  \end{enumerate}
  The norm $\norm{\cdot}$ induces a metric $d(x,y) := \norm{x - y}$ on $V$.
\end{definition}

\begin{definition}[Banach space]
  If $(V,d)$ is a complete normed space,
  we say that $(V,\norm{\cdot})$ is a Banach space.
\end{definition}

\begin{example}[Hilbert space]
  Every Hilbert space is a Banach space where the norm is induced by
  an inner product.
\end{example}

\begin{example}[$\ell_p$ spaces]
  Fix $1 \le p < \infty$ and a natural number $n \geq 1$.
  The space $\ell_p^n$ is the space $\mathbb F^n$ with the $p$-norm:
  \[
    \norm{x}_p :=
    \norm{(x_1,\dots,x_n)}_p =
    \del{\sum_{i=1}^{n} \abs{x_i}^p}^{1/p}.
  \]
  Similarly, the space $\ell^p$ is defined by:
  \[
    \ell_p :=
    \set{(x_n)_{n=1}^{\infty} \,\middle\vert\, \sum_{i=1}^{\infty} \abs{x_i}^p < \infty},
  \]
  with the norm:
  \[
    \norm{x}_p := \del{\sum_{i=1}^{\infty} \abs{x_i}^p}^{1/p}.
  \]
  Finally, we define the space:
  \[
    \ell_{\infty} :=
    \set{(x_n)_{n=1}^{\infty} \,\middle\vert\, \sup_{n} \abs{x_n} < \infty},
  \]
  with the norm $\norm{x}_{\infty} := \sup_{n} \abs{x_n}$.
\end{example}

\begin{remark}
  The space $\ell^2$ is a Hilbert space,
  where the norm is induced by the inner product
  $\ip{x}{y} := \sum_{n=1}^{\infty} x_n \overline{y_n}$.
  This is not true in general for $\ell_p$ spaces.
\end{remark}

\begin{remark}
  For all $1 \le p < \infty$ the space $\ell_p$ is seperable,
  however $\ell_{\infty}$ is a ``very big'' space,
  and it is not seperable.
\end{remark}

\begin{example}[$C$ and $C_0$]
  Define:
  \[
    C :=
    \set{(x_n)_{n=1}^{\infty} \,\middle\vert\, \lim_{n \to \infty} x_n \text{ exists}},
  \]
  with the norm $\norm{x} = \sup_{n} \abs{x_n}$.
  The space $C$ is a subspace of $\ell_{\infty}$,
  and in particular, it is seperable.
  
  Next define:
  \[
    C_0 :=
    \set{(x_n)_{n=1}^{\infty} \,\middle\vert\, \lim_{n \to \infty} x_n = 0}
    \subseteq C \subseteq \ell_{\infty}.
  \]
\end{example}

\begin{remark}
  A subset of a Banach space is a Banach subspace if and only if
  it is closed.
\end{remark}

\begin{example}[The space $C(X)$]
  Let $X$ be a compact topological space.
  We define the space:
  \[
    C(X) :=
    \set{f \colon X \to \mathbb F \mid f \text{ is continuous}},
  \]
  and the norm is $\norm{f} := \norm{f}_{\infty} = \max_{x \in X} \abs{f(x)}$.
\end{example}

\begin{question}
  What does it mean that $f_n \taking{n \to \infty} f$ in
  the Banach space $C(X)$?
\end{question}
\begin{answer}
  It means that $\norm{f_n - f} \taking{n \to \infty} 0$,
  or similarly $\max\abs{f_n - f} \taking{n \to \infty} 0$,
  which is equivalent to the condition $f_n \taking{n \to \infty} f$
  uniformly.
\end{answer}

\begin{example}[Nonexample of a Banach space]
  Define the vector space:
  \[
    C_C(\R^d) :=
    \set{f \colon \R^d \to \R \mid f \text{is cont. and compactly supported}},
  \]
  and consider the norm $\norm{f} := \max\abs{f}$.
  This is a normed space that is \textit{not} complete,
  and therefore this is not a Banach space.

  For example, suppose that $d = 1$ and take $g(x) = e^{-x^2}$.
  % add graph?
  Define the function sequence $\set{f_n}_{n=1}^{\infty}$
  by:
  \[
    f_{n}(x) :=
    \begin{cases}
      e^{-x^2} &x \in [-n,n] \\
      0 &x \notin (-n-1,n+1) \\
      \text{linear} &\text{otherwise}
    \end{cases}.
  \]
  Clearly $\set{f_n}_{n=1}^{\infty} \subseteq C_C(\R)$,
  and $\norm{f_n - g}_{\infty} \taking{n \to \infty} 0$,
  so since it converges, it is a Cauchy sequence,
  however, the limit $g$ is not in $C_C(\R)$,
  which implies that $C_C(\R)$ is not a Banach space.
\end{example}

\begin{remark}
  Remember that every noncomplete metric space is dense in some other
  space and can be completed.
  It turns out that the completion of $C_C(\R^d)$
  is the space
  \[
    C_0(\R^d) :=
    \set{f \colon \R^d \to \R \,\middle\vert\, f \text{ is cont.} \stand
      \lim_{\norm{x} \to \infty} f(x) = 0},
  \]
  with the norm $\norm{f} := \max\abs{f}$.
\end{remark}

% \begin{remark}[A weird sequence]
%   Consider the space of functions %\frac{1}{x^{1/n}}
% \end{remark}

\begin{definition}[$\L^p$ spaces]
  \label{def:Lp-spaces}
  Fix a measure space $(X,\mathcal F,\mu)$
  and a constant $1 \le p < \infty$.
  Define:
  \[
    \L^p(X,\mathcal F,\mu) :=
    \set{f \colon X \to X \,\middle\vert\, f \text{ measurable} \stand
      \del{\int_{X} \abs{f}^p}^{1/p} < \infty} / \sim
  \]
  with the norm $\norm{f}_{p} := \del{\int_{X} \abs{f}^p}^{1/p}$,
  where functions are identified by $f \sim g$ if and only if
  $f =_{\text{a.e.}} g$ (equal almost everywhere).
\end{definition}
\begin{remark}
  Without the equivalence reltaion in \Cref{def:Lp-spaces}
  we could have functions which are not the zero function
  with norm $0$.
\end{remark}

\begin{remark}
  If $X = \N$ and the measure is the counting measure $\mu(A) = |A|$,
  then $\L^p(X,\mu)$ is the previously defined $\ell_p$ space.
\end{remark}

\begin{question}
  We can convince ourself that $\L^p$ spaces are normed spaces,
  but why are they complete?
\end{question}

\begin{exercise}
  A normed space is complete if and only if every absolutely converging
  series converges conditionally.
\end{exercise}

\begin{theorem}
  The space $\L^p(X,\mu)$ is complete.
\end{theorem}
\begin{proof}
  Fix $\set{f_n}_{n=1}^{\infty} \subseteq \L^p(\mu)$ such that
  $\sum_{n=1}^{\infty} \norm{f_n} < \infty$.

  First assume that $\set{f_n}_{n=1}^{\infty}$
  is a sequence of nonnegative functions.
  Define the function $g(x) := \sum_{n=1}^{\infty} f_n(x)$,
  and the sequence $g_N(x) := \sum_{n=1}^{\infty} f_n(x)$.
  Note that
  \[
    \norm{g_N}_{\L^p} \le
    \sum_{n=1}^{N} \norm{f_n} \le
    \underbrace{\sum_{n=1}^{\infty} \norm{f_n}}_{:=\, S} < \infty.
  \]
  Note that $g_N \nearrow g$ pointwise almost everywhere.
  Hence, by MCT we get that
  \[
    \int g^p \dmu =
    \lim_{n \to \infty} \int g_n^p \dmu =
    \lim_{n \to \infty} \norm{g_n}_p^p \le
    S^p < \infty.
  \]
  Therefore $g \in \L^p(\mu)$ and in particular $g < \infty$
  almost everywhere.
  Now check using MCT/DCT that $\norm{g_N - g}_p \taking{N \to \infty} 0$.

  For a general $f_n$ we can decompse it to its positive and negative
  parts:
  \[
    f_n = \alpha_N + i \beta_n =
    \del{\alpha_n^+ - \alpha_n^-} + i(\beta_n^+ + \beta_n^-),
  \]
  we take $\abs{\alpha_n^{\pm}},\abs{\alpha_n^{\pm}} \le \abs{f_n}$,
  so they also satisfy that they are nonpositive/nonnegative
  and satisfy
  \[
    \sum_{n=1}^{\infty} \abs{\alpha_n^{\pm}} < \infty
    \tand
    \sum_{n=1}^{\infty} \abs{\beta_n^{\pm}} < \infty,
  \]
  so $\sum_{n=1}^{\infty} f_n$ is a sum of four convergent sequences in
  and thus converges (in $\L^p$).
\end{proof}

\begin{definition}[The essential supremum]
  Let $(X,\mathcal F,\mu)$ be a measure space,
  and let $g$ be a measurable function.
  We define:
  \[
    \esssup g := \inf \set{m \,\middle\vert\, \mu\del{\set{x \mid g(x) > m} = 0}.}
  \]
\end{definition}

\begin{definition}[$\L^{\infty}$]
  We define the set $\L^{\infty}$ to be the space:
  \[
    \L^{\infty}(X,\mathcal F,\mu) =
    \set{f \colon X \to \mathbb F \mid
      f\text{ measurable and } \esssup|f| < \infty},
  \]
  with the norm $\norm{f}_{\infty} = \esssup |f|$.
\end{definition}

\begin{exercise}
  Let $f \in \L^{p} \cap \L^{\infty}$ for some $1 \le p < \infty$.
  Then $\norm{f}_{\infty} = \lim_{p \to \infty} \norm{f}_p$.
\end{exercise}

\subsection{Linear operators}

\begin{definition}
  Let $X,Y$ be normed spaces over a field $\mathbb F$.
  A linear map $T \colon X \to Y$ is called bounded
  if
  \[
    \norm{T} =
    \sup_{0 \neq x \in X} \frac{\norm{Tx}_Y}{\norm{x}_X} < \infty.
  \]
\end{definition}
\begin{remark}
  We denote the space of bounded operators by $B(X,Y)$.
\end{remark}
\begin{remark}
  Recall that the following conditions
  for are equivalent for a linear operator $T$:
  \begin{enumerate}
    \item[(1)] $T$ is bounded.
    \item[(2)] $T$ is continuous.
    \item[(3)] $T$ is continuous at a point.
    \item[(4)] $T$ is uniformly continuous.
    \item[(5)] $T$ is Lipschitz continuous.
  \end{enumerate}
\end{remark}

\begin{example}[Linear noncontinuous operator]
  Let $\set{v_i}_{i \in I \subset \R}$ be a noncountable basis of $\ell_2$
  (in the sense of linear algebra), which exists by Zorn's lemma.
  We can define $T \colon \ell_2 \to \ell_2$ be the linear map
  $T(v_i) = |i|^{1000} \cdot v_i$.
\end{example}

\begin{proposition}
  Let $Y$ be a Banach space.
  Then $B(X,Y)$ is a Banach space.
\end{proposition}

\begin{definition}[Dual space]
  Let $X$ be a Banach space.
  We define the dual space of $X$ by:
  \[
    X^* := B(X,\mathbb F),
  \]
  which is Banach from the previous theorem.
\end{definition}

\begin{theorem}[Riesz's representation theorem]
  Let $H$ be a Hilbert space.
  Then every $f \in H^*$ is of the form $f_y(x) = \ip{x}{y}$
  for some $y \in H$ and $\norm{y}_{H} = \norm{f_y}_{H^*}$.
\end{theorem}

\begin{remark}
  We have that $H^* \cong H$ by the map $y \mapsto f_y$.
\end{remark}
% over R over C subtlety

\begin{corollary}[Radom--Nikodym theorem]
  Let $\mu,\nu$ be (finite) measures on $(X,\mathcal F)$.
  Suppose that $\mu(A) = 0$ implies $\nu(A) = 0$.
  Then there exists $g \in \L^1(\mu)$ such that $g \geq 0$
  and for every $A \in \mathcal F$ we have $\nu(A) = \int_{A} g \dmu$.
\end{corollary}
\begin{proof}
  Let $H = \L^2(\mu + \nu)$.
  Define the map $l \colon H \to \R$ by $l(f) = \int_X f \dmu$.

  We see that $l$ is bounded by the Cauchy--Schwartz inequality:
  \begin{align*}
    \abs{l(f)} &\le
    \int_X 1 \cdot |f| \dmu \le
    \sqrt{\int_X 1^2 \dmu \cdot \int_X |f|^2 \dmu} \le
    \sqrt{\mu(X)} \sqrt{\int |f|^2 \dif(\mu + \nu)} \\ &=
    \sqrt{\mu(Z)} \norm{f}_{\L^2(\mu + \nu)}.
  \end{align*}
  By Riesz's theorem there exists $h \in H$ such that
  for all $f \in H$ we have:
  \begin{align*}
    \int_X f \dmu &=
    l(f) =
    \ip{f}{h} =
    \int_X f h \dif(\mu + \nu) =
    \int_X f (1 - h) \dmu \\ &= \int_X f h \dnu.
  \end{align*}
  Take $A = \set{x \mid h(x) \le 0}$,
  and $f = \ind_A \in H$.
  We know that
  \[
    \mu(A) =
    \int_A \dmu =
    \int_A (1 - h) \dmu =
    \int_A h \dnu 
  \]
\end{proof}

\begin{remark}
  Note that the inner product is simply:
  \begin{align*}
    \ip{f}{g} = \int_X f \overline g \dif(\mu + \nu) =
    \int_X f \overline g \dmu + \int_X f \overline g \dnu,
  \end{align*}
  and $\norm{f}_{h^2} = \sqrt{\ip{f}{f}}$.
\end{remark}

\begin{theorem}
  Fix $1 \le p < \infty$.
  We have that $(\ell_p)^* = \ell_q$ where $\frac{1}{p} + \frac{1}{q} = 1$.
  Additionally, every $f \in (\ell_p)^*$ is of the form:
  \[
    f_c\del{(a_n)_{n=1}^{\infty}} =
    \sum_{n=1}^{\infty} a_n c_n
  \]
  for some sequence $c = (c_n)_{n=1}^{\infty} \in \ell_q$.
  Moreover, $\norm{f_c}_{(\ell_p)^*} = \norm{c}_{\ell_q}$.
\end{theorem}
% dual of l_1 is l_infty but not the other way around
\begin{proof}
  Assume that $1 < p < \infty$.
  For all $c \in \ell_q$ and $a \in \ell_p$ we have
  \[
    \abs{f_c(a)} =
    \abs{\sum_{n=1}^{\infty} a_n c_n} \le
    \sum_{n=1}^{\infty} |a_n| |c_n| \le_{\text{H\"older}}
    \del{\sum_{n=1}^{\infty} |a_n|^{p}}^{1/p}
    \del{\sum_{n=1}^{\infty} |c_n|^{q}}^{1/q} \le
    \norm{c}_{\ell_q} \norm{a}_{\ell_p},
  \]
  so $f_c$ is well-defined, bounded,
  and $\norm{f_c}_{(\ell_p)^*} \le \norm{c}_{\ell_q}$.

  In order to show equality
\end{proof}


%%%%%% lecture 3

\begin{example}
  Let $X$ be a compact Hausdorff topological space (for example $X = [0,1]$).
  Consider the dual space $C(X)^*$.

  Given a finite Borel measure $\mu$ on $X$,
  we can define the linear operator:
  \[
    \Phi_{\mu}(f) = \int f \dmu \in C(X)^*.
  \]
  Since we have
  \[
    \abs{\Phi_{\mu}(f)} =
    \int |f| \dmu \le
    \mu(X) \cdot \norm{f}_{\infty},
  \]
  we get that $\norm{\Phi_{\mu}}_{C(X)^*} \le \mu(X)$ so it is a bounded
  linear functional.
  Actually, there is equality for example if $f \equiv 1$.
  % why equality
\end{example}

\begin{definition}[Complex measure]
  A complex measure on a measurable space $(X,\mathcal F)$ is a function
  $\mu \colon \mathcal F \to \C$ such that
  for all $\set{A_n}_{n=1}^{\infty} \subseteq \mathcal F$ that are pariwise
  disjoint we have
  $\mu(\bigcup_{n=1}^{\infty} A_n) = \sum_{n=1}^{\infty} \mu(A_n)$.
\end{definition}

\begin{remark}[Hahn decomposition theorem]
  Every complex measure can be written as:
  \[
    \mu = (\mu_1 - \mu_2) + i(\mu_3 - \mu_4),
  \]
  where $\mu_1,\mu_2,\mu_3,\mu_4$ are positive measures,
  and then we define:
  \[
    \int f \dmu =
    \int f \dmu_1 -
    \int f \dmu_2 +
    \int f \dmu_3 -
    \int f \dmu_4.
  \]
\end{remark}

\begin{theorem}[Riesz--Markov--Kakutani representation theorem]
  We have that
  \[
    C(X)^* = M(X) = \set{\text{complex Borel measures on $X$}}.
  \]
  Moreover, we have that:
  \[
    \norm{\Phi_{\mu}}_{C(X)^*} =
    \norm{\mu}_{M(X)} = TV(\mu) = \abs{\mu} =
    \sup\set{\sum_{i=1}^{n} \abs{\mu(A_i)} \mid X = \bigsqcup_{i=1}^{n} A_i}.
  \]
\end{theorem}

\begin{example}
  In a finite (or $\sigma$-finite) meausre space $(X,\mathcal F, \mu)$,
  every $\Phi \in \L^p(X)^*$ is of the form:
  \[
    \Phi_g(f) =
    \int_X f g \dmu
  \]
  for $g \in \L^q(X)$
  and $\norm{\Phi}_{\L^p(X)^*} = \norm{g}_{\L^q}$.

  Proof sketch:
  By H\"older's inequality $\Phi_g \in (\L^p)^*$
  and $\norm{\Phi_g} \le \norm{g}_{\L^q}$,
  and like in the case of $\ell_p$, this is actually an equality.

  In order to show that $\Phi_g$ are all the bounded linear functionals
  use Radon--Nikodym (proof omitted in class).

  Equivalently, there is an isometry $\del{\L^p(X)}^* \cong \L^q(X)$
  for $\frac{1}{p} + \frac{1}{q} = 1$.
\end{example}

\begin{definition}[Isometry]
  A bounded linear functional $T \in \L(X,Y)$ is called an isometry,
  if $T$ is a bijection and
  \[
    \norm{Tx}_Y = \norm{x}_X
  \]
  and we write $X \cong Y$.
  % fix B(X) to \L(X)
\end{definition}

\begin{definition}[Isomorphism]
  A map $T \in \L(X,Y)$ is an isomorphism if $T$ is a bijection and
  $T^{-1} \in \L(Y,X)$.
  Equivalently, for every $T \in \L(X,Y)$ we have
  \[
    c\norm{x} \le \norm{Tx} \le C\norm{x},
  \]
  for some constants $c$ and $C$.
\end{definition}

\begin{example}
  Let $X = \ell_2$ be with the usual norm,
  and let $Y = \ell_2$ with the norm
  \[
    \norm{(a_n)}_{Y} = \norm{(a_n)}_{\ell_2} + \abs{a_1}.
  \]
  We get that $X$ and $Y$ are isomorphic.
  It follows that $Y$ is complete (since $X$ is complete).
\end{example}

\section{The Hahn--Banach theorem}

\begin{remark}
  We claimed that $\ell_1 \subsetneq (\ell_{\infty})^*$.
  How can we find all the elements in $(\ell_{\infty})^*$
  that are not in $\ell_1$?
\end{remark}

\begin{definition}[Sublinear function]
  Let $X$ be a vector space over $\R$.
  A function $p \colon X \to \R$ is said to be sublinear
  if $p(x + y) \le p(x) + p(y)$
  and $p(\lambda x) = \lambda p(x)$ for every $x,y \in X$
  and $\lambda \geq 0$.
\end{definition}
\begin{remark}
  Althought sublinear functions look like norms,
  they are certainly very different.
  For example, some sublinear function do not satisfy $p(x) = p(-x)$.
\end{remark}

\begin{theorem}[The Hahn--Banach theorem, version 1]
  Let $X$ be a vector space over $\R$,
  let $p \colon X \to \R$ be a sublinear function,
  let $Y \subseteq X$ be a subspace,
  and let $l \colon Y \to \R$ be a linear function.
  If $l(y) \le p(y)$ for all $y \in Y$,
  then there exists a linear extension $\tilde l \colon X \to \R$
  such that $\tilde l\vert_Y = \ell$ and $\tilde l(x) \le p(x)$ for all $x \in X$.
\end{theorem}
\begin{proof}
  Assume that $X = \vspan\set{Y,x_0}$ for some $x_0 \in X$,
  so that every $x \in X$ can be written as $x = y + \lambda x_0$
  for some $y \in Y$ and $\lambda \in \R$.

  Note that
  \[
    \tilde l(x) =
    \tilde l(y + \lambda x_0) =
    \tilde l(y) + \lambda \tilde l(x_0) =
    \ell(y) + \lambda \tilde l(x_0),
  \]
  so we only get to choose $a = \tilde l(x_0)$.
  Can we choose $a$ such that $l(y) + \lambda a = \tilde l(y + \lambda x_0)
  \le p(y + \lambda x_0)$.

  For $\lambda = 0$ any value of $a$ is fine,
  for $\lambda > 0$ we can solve for $a$ from
  $l(y) + \lambda a = \tilde l(y + \lambda x_0) \le p(y + \lambda x_0)$
  and get that for every $y \in Y$ we need
  $a \le \frac{p(y + \lambda x_0) - l(y)}{\lambda}$,
  and similarly we can find restriction for when $\lambda < 0$,
  and conclude that the restrictions are:
  \[
    \frac{l(z) - p(z + \lambda x_0)}{\lambda}
    \le a \le
    \frac{p(y + \lambda x_0) - l(y)}{\lambda},
  \]
  or equivalently, we need to show that exists $a$ such that
  \[
    \sup_{\mu > 0,z \in Y}\set{l(z/\mu) - p(z/\mu - x_0)}
    \le a \le
    \sup_{\lambda > 0,y \in Y}\set{p(y/\lambda + x_0) - l(y/\lambda)}
  \]
  % picture

  In the general case, if we need the extend $Y$ to $X$ by a finite
  amount of variables, we repeat this process.
  In the infinte case, we need to use Zorn's lemma.
  Consider:
  \[
    \mathcal C = \set{(Z,\tilde l) \mid \tilde l\vert_Y = \ell \stand \tilde l \le p \text{ on }Z},
  \]
  for $Z \subseteq Y \subseteq X$.
  We can order $\mathcal C$ be saying that
  $(Z_1,\tilde l_1) \le (Z_2,\tilde l_2)$ if $Z_1 \subseteq Z_2$
  and $\tilde l_2\vert_{Z_1} = \tilde l_1$.

  Take $\set{(Z_i,\tilde l_i)}_{i \in I}$ to be a chain in $\mathcal C$.
  Take $Z = \bigcup_{i \in I} Z_i$ and $l$ be the ``obvious''
  extension (every $z \in Z$ belongs to some $Z_i$,
  take $l(z) = \tilde l_i(z)$, since they agree by the chain containment).
  Thus, $(Z,l)$ is an upper bound fo the chain.
  By Zorn's lemma $C$ has a maximal element, $(Z, \ell)$.
  By step $1$, $Z = X$.
\end{proof}

\begin{theorem}[The Hahn--Banach theorem, version 2]
  \label{thm:hahn-banach-ver-two}
  Let $X$ be a normed space over (for example a Banahc space),
  let $Y \subseteq X$ be a subspace,
  and let $l \in Y^*$ be a continuous functional.
  Then there exists some linear functional
  $\tilde l \in X^*$ such that $\tilde l \vert_{Y} = l$
  and $\|\tilde l\|_{X^*} = \norm{l}_{Y^*}$.
\end{theorem}

\begin{example}
  Let $\ell_{\infty}$ and $Y = C$.
  Take $f \in Y*$ to be the functional defined by
  \[
    f\del{(a_n)_{n=1}^{\infty}} = \lim_{n \to \infty} a_n.
  \]
  Note that
  \begin{align*}
    \abs{f\del{(a_n)}} =
    \abs{\lim_{n \to \infty} a_n} =
    \lim_{n \to \infty} |a_n| \le \sup_{n} \abs{a_n} = \norm{a}_Y,
  \end{align*}
  it follows that $\norm{f}_{Y^*} \le 1$.
  However, we notice that for $a = (1,1,1,\dots)$ we see
  $\norm{f}_{Y^*} = 1$.

  By the Hahn--Banach theorem,
  we can extend $f$ to $\tilde f \in (\ell_{\infty})^*$
  such that $\|\tilde f\|_{(\ell_{\infty})^*}$.

  Note that $\tilde f$ is not of the form:
  \[
    \tilde f \del{(a_n)_{n=1}^{\infty}} =
    \sum_{n=1}^{\infty} a_n b_n \quad (b_n) \in \ell_1,
  \]
  since otherwise, $b_n = \tilde f(e_n) = f(e_n) = 0$.
  Thus $b_n = 0$ for all $n$ and then $\tilde f \equiv 0$ in contradiction
  to $\norm{f}_{Y^*} = 1$.
\end{example}

\begin{proofof}{\Cref{thm:hahn-banach-ver-two}}
  Choose $p(x) = \norm{f}_{Y^*} \norm{x}_X$.
  Clearly, $p(x)$ is sublinear.
  On $Y$, we have $l(y) \le p(y) = \norm{l}_{Y^*} \norm{y}$.
  By version 1 of the Hahn--Banach theorem,
  there exists an extension $\tilde l$
  of $l$ such that $\tilde l(x) \le p(x) = \norm{\ell}_{Y^*} \norm{x}$
  for all $x \in X$.
  But then also
  \[
    -\tilde l(x) = \tilde l(-x) \le p(-x) = \norm{l}_{Y^*} \norm{x},
  \]
  and thus
  \[
    \abs{\tilde l(x)} \le
    \norm{l}_{Y^*} \norm{x} \implies
    \norm{\tilde l}_{X^*} \le \norm{l}_{Y^*}.
  \]
  The inequality $\|\tilde l\|_{X^*} \geq \norm{l}_{Y^*}$ is true for all
  extensions.
  This completes the proof for real normed spaces.

  For the case of complex normed spaces, we decompose $l$ into its
  real and complex parts, and use some algebra and the proof of the real
  case in order to complete the proof.
\end{proofof}

\begin{corollary}
  \label{cor:hahn-banach-one}
  Let $X$ be a Banach space.
  Then for every $x \in X$ there exists $f \in X^*$ such that
  $\norm{f} = 1$ and $f(x) = \norm{x}$.
\end{corollary}
\begin{proof}
  Let $Y = \vspan{x} \subseteq X$,
  and define $f \colon Y \to \mathbb F$ by $f(\lambda x) = \lambda \norm{x}$.

  We have that $f$ is linear and $\norm{f} = 1$ and
  \[
    \abs{f(\lambda x)} =
    \abs{\lambda} \norm{x} =
    \norm{\lambda x}.
  \]
  Extend $f$ by the Hahn--Banach theorem to conclude the proof.
\end{proof}

\begin{corollary}
  For all $x \neq y \in X$ there exists $f \in X^*$ such that
  $f(x) \neq f(y)$.
\end{corollary}
\begin{proof}
  By \Cref{cor:hahn-banach-one}
  take $f$ with $\norm{f} = 1$
  and then:
  \[
    f(x) - f(y) =
    f(x - y) =
    \norm{x - y} \neq = 0.
  \]
  % go over
\end{proof}

\begin{corollary}
  Let $X$ be a normed space,
  let $E \subseteq X$ be a subspace,
  and let $x_0 \in X$ be a point such that $d(x_0,E) = d > 0$.
  Then there exists $f \in X^*$ such that $\norm{f} = 1$,
  and $f(E) = 0$, and $f(x_0) = d$.
\end{corollary}
\begin{proof}
  Take $Y = \vspan\set{E,x_0}$.
  Define $f \in Y^*$ by $f(e + \lambda x_0) = \lambda d$.

  For all $e \in E$ and $\lambda \in \mathbb F$
  \[
    \norm{e + \lambda x_0} =
    \norm{\lambda\del{x_0 - (- \frac{e}{\lambda})}} =
    \abs{\lambda} \cdot \norm{x_0 - \underline{(- \frac{e}{\lambda})}_{\in E}} \geq
    \abs{\lambda} d =
    \abs{f(e + \lambda x_0)},
  \]
  it follows that $\norm{f}_{Y^*} \le 1$ (in fact, there is an equality).
  We now extend by Hahn--Banach to complete the proof.
\end{proof}



\end{document}
